\documentclass[12pt,letterpaper]{report}
\usepackage[pdftex]{graphicx}
\usepackage[top=1in,bottom=1in,left=1in,right=1in]{geometry}
%\usepackage{fancyhdr}

%\setlength{\headheight}{15pt}
%\pagestyle{fancy}

\newcommand{\HRule}{\rule{\linewidth}{0.5mm}}

\begin{document}
 \begin{titlepage}

\begin{center}

\includegraphics[width=0.5\textwidth]{./resources/logo}\\[1cm]

\textsc{\Huge ASCENDS Data Visualization}\\[1.5cm]

\textsc{\Large DEVELOP, NASA Langley Research Center}\\[0.5cm]

\HRule \\[0.4cm]
{ \huge \bfseries Technical Report \emph{Draft}}\\[0.4cm]

\HRule \\[1.5cm]

\begin{minipage}{0.45\textwidth}
\begin{flushleft} \large
\emph{Principle Investigators:}\\
T. Nelson \textsc{Hillyer}\\
Jacob T.\textsc{Atkins}\\
Benjamin S. \textsc{Hughes}\\
Andrew R. \textsc{Kramolisch}\\
Mark A. \textsc{Lotts}
\end{flushleft}
\end{minipage}
\begin{minipage}{0.45\textwidth}
\begin{flushright} \large
\emph{Science Advisors:}\\
Dr. F. Wallace \textsc{Harrison}\\
Dr. Edward \textsc{Browell}
\end{flushright}
\end{minipage}

\vfill

{\large \today}

\end{center}

\end{titlepage}



 \begin{abstract}
  The Active Sensing of Carbon dioxide Emissions over Nights, Days, and Seasons (ASCENDS) mission is a NASA decadal survey mission that aims to enhance understanding of carbon dioxide ($CO_2$) in the global carbon cycle through three scientific objectives: quantifying global spatial distributions of atmospheric $CO_2$, quantifying current distributions of terrestrial and oceanic sources and sinks of $CO_2$ on $1^\circ$ by $1^\circ$ grids at weekly resolution, and providing a scientific basis for future $CO_2$ projections of sources and sinks through data-driven enhancements of Earth system process modeling. The purpose of this project was to create high resolution visualizations of the ASCENDS mission data to help in the analysis and utilization of the $CO_2$ readings. The ASCENDS mission steering committee at NASA Langley Research Center requested a software product that could dynamically generate KML files for visualization using the data gathered by a series of flights over areas of interest with specially outfitted aircraft which simulated the data that would be retrieved by an ASCENDS satellite. The ASCENDS Data Visualization Team created the requested visualization tool that can be used to display ASCENDS data from future flights. In the future, this tool can be used to assist in the funding of the ASCENDS mission to improve climate and air quality research.
 \end{abstract}

 \section*{Introduction}
  \subsection*{ASCENDS Mission}
   \paragraph{}
    The Active Sensing of Carbon Dioxide Emissions over Nights, Days, and Seasons (ASCENDS) mission is currently an exploratory mission using simultaneous laser remote sensing of $CO_2$ and $O_2$. The technology employed was based on a recommendation from the National Research Council of the United States National Academies. Active sensing allows the ASCENDS team to enhance understanding of the role of $CO_2$ in the global carbon cycle, due to its lack of seasonal, latitudinal, or diurnal bias.
  \subsection*{Visualization}
   \paragraph{}
    The ASCENDS Team project involved visualizing the data gathered with active sensing technology. Whereas a barrage of numbers is overwhelming, visualizations allow for clarity with a basic understanding of the data. Additionally, visualizations can display an enormous amount of data in an easily accessible format. Because of the technology behind ASCENDS, the visuals generated by the ASCENDS Team project can be compared and contrasted with other $CO_2$ models accurately.
   \paragraph{}
    Generation of the visualizations is accomplished by first piecing the data together in a relational database using the Gigantron framework. Gigantron is an in-house framework designed to ease the process of importing data, linking it together, and manipulating it. Gigantron runs on top of JRuby, a Ruby implementation on the Java Virtual Machine, for speed and rapid development. Rake tasks inside the generated Gigantron project allow for the dynamic generation of Keyhole Markup Language (KML) files, by utilizing ActiveRecord and ruby-kml.
  \subsection*{Importance}
   \paragraph{}
    The ability of visualizations to be utilized and promote a better understanding of $CO_2$ movement patterns explicates the necessity for an ASCENDS project satellite. Having more data recorded and imported into the visualization will eventually allow for a global analysis of atmospheric $CO_2$ levels. The data retrieved from the HYSPLIT model allows for the visualization of the back trajectories of carbon dioxide plumes, thus allowing for the study of the movement and sources of carbon dioxide. In conjunction with the in-situ data, the individual time trial, or ITT, data showed the carbon dioxide mixing ratio with the wind. The differences between the observed data and the predicted data permitted a high-resolution analysis of the carbon dioxide levels. The predictive models closely resembled the observed data, implying that the US will be a net source of $CO_2$.
 \section*{Methodology}
  \subsection*{Data Obtainment}
   \paragraph{}
    The first step in the project was to obtain raw data from the ASCENDS Team's various flights in a plane outfitted with $CO_2$ sensing equipment over eastern Virginia and North Carolina.  The data was provided in multiple spreadsheets and raw text files, and included information about readings such as wind speed, direction, latitude, longitude, $CO_2$ concentrations, and date. Time stamps connected the data across multiple files. One set of data from the ITT information contained the average $CO_2$ concentration over the column of air below the aircraft. The second set of data was the in-situ data, which was the mixing ratio of $CO_2$ as compared to the rest of the air.
  \subsection*{Data Visualization}
   \paragraph{}
    Once the data files were obtained, scripts written in JRuby analyze and parse essential information. The resulting data was then written into a KML file using the ruby-kml library. The algorithms in the program were used to display the ITT data in column format, with the in-situ data on top of the ITT columns. A separate algorithm was used to color-code the resulting polygonal columns according to the level of $CO_2$ in the air, using the established ASCENDS color bar. The resulting KML file was imported into Google Earth, where the latitude, longitude, and altitude data were used to map out the columns on the exact place where the data points were gathered. This allows the data to be analyzed in both a geographical and visual sense.
  \subsection*{Model Visualization and Comparison}
   \paragraph{}
  Data from various predictive models was written into KML files and imported into Google Earth to be compared with the observed data. A cross-sectional analysis of the observed data versus the predictive data versus the predictive data was visualized based on both the time and location of data point. The reason for this is to dispel anomalies that may be caused by the dynamic system from which the observational data was taken.  
 \section*{Results}
  \paragraph{}
The result of the project is primarily the Google Earth visualization of the carbon dioxide data. The in-situ data and the ITT data are displayed in columns along the flight path of the plane. By default, the HYSPLIT wind vectors are displayed every sixty seconds along the plane's route, but that interval can be changed as needed. Finally, specific analysis of portions of the data work by breaking the flight into minutes, with every data point for that minute and other manually selected minutes available for individual study. This can also display future mission data and will continue to be useful for the ASCENDS mission.

 \section*{Discussion}
  \paragraph{}
   Carbon dioxide in the atmosphere can be visualized according to the exact geographic location. This allows both scientists and politicians to better interpret the visualized data. The display of in-situ mixing ratio data alongside the ITT data allows for a more in-depth analysis of how carbon dioxide levels affect the composition of the air in specific atmospheric locations. This ASCENDS visualization will be useful in helping to calibrate future visualizations and can serve as a benchmark for $CO_2$ analysis in future projects.
 \section*{Conclusion}
  \paragraph{}
   As global warming and carbon dioxide management directly affect the future of the Earth, air pollution is becoming an increasingly relevant topic.  This project is relevant because it allows for the advancement of environmental research as well as hopefully leading to scientific breakthroughs that will help the environment's long-term sustainability.
  \paragraph{}
  High resolution carbon dioxide research will now be possible, thus helping in the fight against pollution and global warming. An ASCENDS satellite will allow for global carbon dioxide atmospheric studies. The visualization could be improved by adding more carbon dioxide predictive models. This will allow so that the observed data can be compared to multiple expected values, thus obtaining a more accurate result. Optimizing the visualization algorithms will improve the product and make it less computationally demanding. A friendlier user interface would make the visualization more accessible to people who have minimal experience with Google Earth and KML files. This visualization must be used by the scientific community to aid research and eventually discover solutions to the world's environmental problems. The visualization will ultimately help the ASCENDS mission to gain more attention.
 \section*{Acknowledgments}

\end{document}


